\documentclass[11pt,a4paper]{article}

\usepackage{fontspec}
\IfFontExistsTF{Times New Roman}{\setmainfont{Times New Roman}}{\setmainfont{Tinos}}
\IfFontExistsTF{Consolas}{\setmonofont{Consolas}[Scale=0.92]}{\setmonofont{Liberation Mono}[Scale=0.92]}
\usepackage[english]{babel}
\usepackage{geometry}
\usepackage{xcolor}
\usepackage{enumitem}
\usepackage{listings}
\usepackage{booktabs}
\usepackage{array}
\usepackage{amsmath,amssymb}
\usepackage{titlesec}
\usepackage{fancyhdr}
\usepackage{parskip}
\usepackage{tcolorbox}
\tcbuselibrary{skins}
\usepackage{tabularx}
\usepackage{lastpage}
\usepackage{hyperref}

\geometry{top=1.35cm,bottom=1.35cm,left=1.65cm,right=1.65cm}

\definecolor{primaryblue}{RGB}{18,65,115}
\definecolor{secondaryteal}{RGB}{14,102,110}
\definecolor{SoftBlue}{RGB}{232,242,250}
\definecolor{CodeBg}{RGB}{248,250,252}
\definecolor{DarkGray}{RGB}{55,55,55}
\definecolor{boxbg}{RGB}{248,250,252}
\definecolor{borderblue}{RGB}{150,185,215}

\setlength{\parindent}{0pt}
\setlength{\parskip}{0pt}
\setlist[itemize]{leftmargin=1.3em,itemsep=1pt,topsep=1pt,parsep=0pt,partopsep=0pt}
\setlist[enumerate]{leftmargin=1.4em,itemsep=1pt,topsep=1pt,parsep=0pt,partopsep=0pt}
\setlist[enumerate,2]{topsep=0.5pt,itemsep=0pt,parsep=0pt,partopsep=0pt}

\lstset{
  language=Python,
  basicstyle=\ttfamily\scriptsize,
  keywordstyle=\color{blue!75!black}\bfseries,
  stringstyle=\color{red!65!black},
  commentstyle=\color{primaryblue!80!black}\itshape,
  showstringspaces=false,
  columns=fullflexible,
  keepspaces=true,
  breaklines=true,
  frame=single,
  rulecolor=\color{primaryblue!45},
  backgroundcolor=\color{CodeBg},
  xleftmargin=0.2em,
  xrightmargin=0.2em,
  aboveskip=0.2em,
  belowskip=0.2em,
  tabsize=4
}

\pagestyle{fancy}
\fancyhf{}
\setlength{\headheight}{14pt}
\fancyhead[L]{\footnotesize\textbf{DSAI1003 -- Fundamental Programming Concepts in Python}}
\fancyhead[R]{\footnotesize\textbf{Functions \& Modular Design Exam -- 60 Minutes}}
\fancyfoot[L]{\scriptsize Faculty of Data Science and Artificial Intelligence -- National Economics University}
\fancyfoot[R]{\footnotesize Page \thepage\ of \pageref{LastPage}}
\renewcommand{\headrulewidth}{0.4pt}
\renewcommand{\footrulewidth}{0.3pt}

\titleformat{\section}{\large\bfseries\color{primaryblue}}{}{0pt}{}
\titleformat{\subsection}{\normalsize\bfseries\color{primaryblue}}{}{0pt}{}
\titlespacing*{\section}{0pt}{3pt}{1pt}
\titlespacing*{\subsection}{0pt}{2pt}{1pt}

\newtcolorbox{headerbox}{
  colback=white,
  colframe=primaryblue,
  boxrule=0.7pt,
  arc=1.5mm,
  left=2.5mm,
  right=2.5mm,
  top=1.0mm,
  bottom=1.0mm
}

\newtcolorbox{answerbox}[2][]{
  enhanced,
  colback=white,
  colframe=DarkGray!45,
  boxrule=0.5pt,
  arc=1mm,
  left=2mm,
  right=2mm,
  top=1.0mm,
  bottom=1.0mm,
  title=\footnotesize\color{DarkGray}\textbf{#2},
  #1
}

\hypersetup{
  colorlinks=true,
  linkcolor=primaryblue,
  urlcolor=primaryblue,
  pdftitle={DSAI1003 - Functions and Modular Design in Python Exam},
  pdfauthor={Faculty of Data Science and Artificial Intelligence - National Economics University (NEU)}
}

\begin{document}

% =========================================================================
% PAGE 1: OFFICIAL HEADER + ANSWER SHEET + SECTION A (Q1 - Q5)
% =========================================================================

\noindent
\begin{minipage}[t]{0.64\textwidth}
  \raggedright\small
  \textbf{NATIONAL ECONOMICS UNIVERSITY}\\
  \textbf{COLLEGE OF TECHNOLOGY}\\
  \textbf{FACULTY OF DATA SCIENCE \& ARTIFICIAL INTELLIGENCE (FDA)}
\end{minipage}%
\hfill
\begin{minipage}[t]{0.34\textwidth}
  \raggedleft\footnotesize
  \textbf{OFFICIAL PROGRESS EXAMINATION}\\
  Academic Year: 2026--2027\\
  Semester: Fall 2026 \quad Course Code: \textbf{DSAI1003}
\end{minipage}

\vspace{0.1em}

\begin{center}
  {\Large\bfseries\color{primaryblue} FUNDAMENTAL PROGRAMMING CONCEPTS IN PYTHON}\\[0.05em]
  {\normalsize\bfseries PROGRESS EXAM: FUNCTIONS, MODULAR DESIGN \& VARIABLE SCOPE}\\[0.05em]
  {\normalsize\textbf{Paper Code:} \rule{1.6cm}{0.4pt} \quad | \quad \textbf{Duration: 60 Minutes} (Excluding paper distribution)}
\end{center}

\vspace{0.05em}

\begin{headerbox}
\begin{tabular*}{\textwidth}{@{\extracolsep{\fill}}ll@{}}
  \textbf{Candidate Full Name:} \dotfill & \textbf{Student ID (MSSV):} \dotfill \\[0.15em]
  \textbf{Class / Cohort:} \dotfill & \textbf{Exam Room / Seat No.:} \dotfill
\end{tabular*}

\vspace{0.15em}
\renewcommand{\arraystretch}{1.0}
\setlength{\tabcolsep}{4.5pt}
\centering
\begin{tabular}{|l|c|c|c|c||c|c|}
\hline
\textbf{Section} & \textbf{A} & \textbf{B} & \textbf{C} & \textbf{D} & \textbf{Total (100)} & \textbf{Examiner Signature} \\
\hline
\textbf{Max Score} & 20 & 25 & 25 & 30 & \textbf{100} & \\[0.25em]
\cline{1-6}
\textbf{Score} & & & & & & \\[0.35em]
\hline
\end{tabular}

\vspace{0.1em}
\raggedright
\scriptsize
\textbf{Instructions:} Closed-book examination. All electronic devices and internet access are prohibited. Answer all questions directly in the designated spaces. In programming tasks, apply modular programming principles, function definitions (\texttt{def}), parameter passing, return statements, and proper variable scope. Do \textbf{not} use global variables or third-party libraries. Organize scripts cleanly with helper functions and a designated \texttt{main()} driver. Do not detach pages.
\end{headerbox}

\vspace{0.5mm}
\noindent
\textbf{MULTIPLE-CHOICE ANSWER SHEET} \textit{(Write A, B, C, or D in the corresponding box)}:
\begin{center}
\setlength{\tabcolsep}{5pt}
\renewcommand{\arraystretch}{1.0}
\begin{tabular}{|c|c|c|c|c|c|c|c|c|c|c|}
\hline
\textbf{Question} & \textbf{1} & \textbf{2} & \textbf{3} & \textbf{4} & \textbf{5} & \textbf{6} & \textbf{7} & \textbf{8} & \textbf{9} & \textbf{10} \\
\hline
\textbf{Answer} & & & & & & & & & & \\[1.5mm]
\hline
\end{tabular}
\end{center}
\vspace{-3mm}
\section*{Section A: Functions, Modular Design \& Scope Concepts \hfill \normalsize\normalfont(20 points)}
\textit{Answer all 10 questions. Each correct answer carries 2 points. Suggested time: 12 minutes.}

\begin{enumerate}[label=\textbf{Q\arabic*.},leftmargin=2em,itemsep=1pt,topsep=1pt]

\item What is the central engineering principle behind treating a function as a ``black box''?
\begin{enumerate}[label=\Alph*.,leftmargin=2em]
  \item The caller must inspect every internal line of code before making a function call.
  \item The caller relies only on the function's interface (inputs, outputs, specification) without needing internal details.
  \item A black-box function cannot accept parameters or return computed values to the caller.
  \item It guarantees that the function executes with zero memory consumption.
\end{enumerate}

\item What is the precise distinction between a \textbf{parameter} and an \textbf{argument} in Python?
\begin{enumerate}[label=\Alph*.,leftmargin=2em]
  \item A parameter is the passed value; an argument is the variable declared in the function header.
  \item Parameters and arguments are identical synonyms with no syntactic or conceptual distinction.
  \item A parameter is the variable in the function header; an argument is the concrete value supplied in the call.
  \item Arguments can only be string literals; parameters can only be numeric types.
\end{enumerate}

\item What is the fundamental operational difference between \texttt{return value} and \texttt{print(value)}?
\begin{enumerate}[label=\Alph*.,leftmargin=2em]
  \item \texttt{print()} passes data back to the calling expression; \texttt{return} only displays characters on screen.
  \item \texttt{return} sends computed data back to the caller for further use; \texttt{print()} displays text and returns \texttt{None}.
  \item \texttt{print()} halts the function immediately; \texttt{return} allows subsequent lines to execute.
  \item Both statements are functionally identical and interchangeable in arithmetic expressions.
\end{enumerate}

\item In Python, if a function finishes executing its body without a \texttt{return} statement, what does it return?
\begin{enumerate}[label=\Alph*.,leftmargin=2em]
  \item The integer \texttt{0}.
  \item An empty string \texttt{""}.
  \item The special singleton value \texttt{None} of type \texttt{NoneType}.
  \item A fatal compile-time \texttt{MissingReturnError} is raised.
\end{enumerate}

\item Given \texttt{def modify(x): x = x + 10}, if a caller runs \texttt{a = 5} followed by \texttt{modify(a)}, what is \texttt{a}?
\begin{enumerate}[label=\Alph*.,leftmargin=2em]
  \item \texttt{15}, because integer arguments are passed by reference and modified in place.
  \item \texttt{5}, because parameter \texttt{x} is a local variable; reassigning \texttt{x} does not modify caller variable \texttt{a}.
  \item \texttt{10}, because the parameter replaces the caller's variable in memory.
  \item An \texttt{UnboundLocalError} occurs because \texttt{a} was not declared as global.
\end{enumerate}

\end{enumerate}

\newpage

% ==============================================================================
% PAGE 2: SECTION A (Q6 - Q10) + SECTION B Q11
% ==============================================================================

\begin{enumerate}[label=\textbf{Q\arabic*.},leftmargin=2em,itemsep=1.5pt,topsep=1pt,resume]

\item Which statement correctly defines the \textbf{lifetime} and \textbf{scope} of a local variable inside a function?
\begin{enumerate}[label=\Alph*.,leftmargin=2em]
  \item It exists for the entire program duration and can be accessed from any function.
  \item It is created when the function executes and destroyed upon exit; it cannot be referenced outside.
  \item It retains its value across separate function calls without using the \texttt{global} keyword.
  \item It is saved directly to disk and reloaded whenever any module is imported.
\end{enumerate}

\item What happens when a local variable inside a function shares the same name as an existing global variable?
\begin{enumerate}[label=\Alph*.,leftmargin=2em]
  \item Python raises a \texttt{NameCollisionError} when the script starts.
  \item The global variable overrides the local variable, preventing the local variable from being assigned.
  \item The local variable \textbf{shadows} (hides) the global variable within the function's local scope.
  \item Python automatically merges both variables into a two-element tuple.
\end{enumerate}

\item In the engineering process of \textbf{stepwise refinement} (top-down design), how should a complex problem be approached?
\begin{enumerate}[label=\Alph*.,leftmargin=2em]
  \item Write all calculations sequentially inside one continuous global script without functions.
  \item Decompose the complex problem into smaller sub-tasks, each implemented as an independent helper function.
  \item Always implement recursion first before considering iterative or procedural decomposition.
  \item Combine user inputs, validations, mathematical formulas, and screen formatting into a single function.
\end{enumerate}

\item In software development, what is a function \textbf{stub}?
\begin{enumerate}[label=\Alph*.,leftmargin=2em]
  \item A permanent syntax error that halts Python compilation.
  \item A temporary placeholder implementation with a valid signature that returns a default value for incremental testing.
  \item A built-in Python keyword used to terminate infinite loops.
  \item A compiled binary extension written in C.
\end{enumerate}

\item Why is organizing program logic within a \texttt{def main():} function considered a best practice in Python?
\begin{enumerate}[label=\Alph*.,leftmargin=2em]
  \item It accelerates Python bytecode execution by 100 times.
  \item It is mandatory because Python cannot run any statements outside a function named \texttt{main}.
  \item It prevents accidental global variables, provides a clear program entry point, and promotes modular testing.
  \item It automatically converts all input variables into floating-point numbers.
\end{enumerate}

\end{enumerate}

\section*{Section B: Output Tracing, Scope Diagnosis \& Error Correction \hfill \normalsize\normalfont(25 points)}
\textit{Answer all 5 questions. Each question carries 5 points. Suggested time: 15 minutes.}

\subsection*{Q11. Exact Output Tracing with Parameter Passing \& Return Tuples (5 points)}
Predict the \textbf{exact screen output}, including punctuation, spaces, and line breaks:

\begin{lstlisting}
def compute_discount(subtotal, rate=0.10):
    discount = subtotal * rate
    return discount

def finalize_bill(amount, is_member):
    tax = amount * 0.08
    rate = 0.15 if is_member else 0.05
    disc = compute_discount(amount, rate)
    net_total = amount - disc + tax
    return disc, net_total

total = 100.0
discount, final_due = finalize_bill(total, True)
print(f"Subtotal: ${total:.2f} | Discount: ${discount:.2f}")
print(f"Final Due: ${final_due:.2f}")
print(f"Saved: {discount >= 15.0 and total == 100.0}")
\end{lstlisting}

\begin{answerbox}[height=2.8cm]{Output Box for Q11}
\end{answerbox}

\newpage

% ==============================================================================
% PAGE 3: SECTION B, Q12--Q15
% ==============================================================================

\subsection*{Q12. Evaluation of Function Calls and Return Values (5 points)}
Write the exact evaluated value (or boolean result \texttt{True} / \texttt{False}) for each expression:

\begin{enumerate}[label=\alph*),leftmargin=2em,itemsep=1pt]
  \item \texttt{def notify(msg): print(msg)}\\
  Result of \texttt{notify("Done") is None}: \hrulefill
  \item \texttt{def cap(v, m=50): return m if v > m else v * 2}\\
  Result of \texttt{cap(20) + cap(60)}: \hrulefill
  \item \texttt{def valid(c): return len(c) == 6 and c[:2].isalpha() and c[2:].isdigit()}\\
  Result of \texttt{valid("NE2026")}: \hrulefill
  \item \texttt{def step(x): x = x + 5; return x * 2} with \texttt{a = 4; b = step(a)}\\
  Value of \texttt{a + b}: \hrulefill
  \item Evaluated result of \texttt{max(round(3.456, 1), min(4.2, 3.8))}: \hrulefill
\end{enumerate}

\subsection*{Q13. Diagnose and Correct Common Function Defects (5 points)}
For each code scenario, state the underlying defect and provide the corrected Python statement.

\renewcommand{\arraystretch}{1.1}
\begin{tabularx}{\textwidth}{|p{5.1cm}|X|X|}
\hline
\textbf{Code / Situation} & \textbf{Problem / Logic Flaw} & \textbf{Correction / Valid Code} \\
\hline
\texttt{def calc\_area(w, h):}\\
\texttt{\ \ \ \ print(w * h)}\\
\texttt{ans = calc\_area(5, 4) + 10} & & \\[2.0mm]
\hline
\texttt{def reset\_counter(c):}\\
\texttt{\ \ \ \ c = 0}\\
\texttt{count = 15}\\
\texttt{reset\_counter(count)} & & \\[2.0mm]
\hline
\texttt{def prepare\_total():}\\
\texttt{\ \ \ \ subtotal = 100.0}\\
\texttt{prepare\_total()}\\
\texttt{print(subtotal)} & & \\[2.0mm]
\hline
\end{tabularx}

\subsection*{Q14. Variable Scope, Name Shadowing \& State Isolation (5 points)}
Predict the exact screen output of the following script:

\begin{lstlisting}
rate = 0.05

def apply_surcharge(base):
    rate = 0.12
    surcharge = base * rate
    return surcharge

def process_order(base):
    fee = apply_surcharge(base)
    total = base + fee + (base * rate)
    return total

amt = 200.0
due = process_order(amt)
print(f"Global Rate: {rate * 100:.0f}%")
print(f"Payable Due: ${due:.2f}")
print(f"Base Unchanged: {amt == 200.0}")
\end{lstlisting}

\begin{answerbox}[height=1.5cm]{Output Box for Q14}
\end{answerbox}

\subsection*{Q15. Function Specification, Docstrings \& Stubs (5 points)}
\begin{itemize}[leftmargin=2em]
  \item Write a standard Python \textbf{docstring} for a function \texttt{calculate\_fuel(distance, consumption\_rate)} explaining parameters, pre-conditions, and return value:\\
  \rule{\linewidth}{0.4pt}
  \item Write a syntactically valid \textbf{function stub} for \texttt{verify\_token(token\_str)} that accepts a string and returns a safe default boolean value \texttt{False}:\\
  \rule{\linewidth}{0.4pt}
\end{itemize}

\newpage

% ==============================================================================
% PAGE 4: SECTION C
% ==============================================================================

\section*{Section C: Program Design \& Modular Decomposition \hfill \normalsize\normalfont(25 points)}
\textit{Answer all tasks. Suggested time: 15 minutes.}

\begin{tcolorbox}[colback=boxbg,colframe=secondaryteal,arc=1.5mm,boxrule=0.6pt,top=1.5mm,bottom=1.5mm]
\textbf{Problem Specification: Logistics Fleet Trip Reimbursement Engine}\\[0.5mm]
\small
A logistics dispatch platform calculates driver trip compensation based on travel distance $d$ (km), cargo weight $w$ (kg), driving duration $t$ (hours), shift timing (\texttt{is\_night}), and diesel fuel price per liter $p$:
\begin{itemize}[itemsep=0.5pt,topsep=0.5pt]
  \item \textbf{Input Validation:} Distance must satisfy $d > 0$; cargo weight must satisfy $0 < w \le 5000.0$; driving time must satisfy $t > 0$. Otherwise, reject as invalid.
  \item \textbf{Fuel Cost Computation:} Base fuel consumption is 10.0~L/100km. If cargo weight exceeds 2000.0~kg, heavy-load friction increases consumption by 15\% (i.e. 11.5~L/100km). Fuel cost is: $\text{liters} \times p$, where $\text{liters} = (d \times \text{consumption\_rate}) / 100.0$.
  \item \textbf{Driver Pay Calculation:} Driver receives a base wage of \$20.00/hour plus a mileage allowance of \$0.30/km. If \texttt{is\_night} is \texttt{True}, the hourly rate receives a 25\% premium (\$25.00/hour).
  \item \textbf{Total Reimbursement:} Sum of fuel cost and driver pay. If total trip distance exceeds 400.0~km, add a flat \$50.00 long-haul bonus.
\end{itemize}
\end{tcolorbox}

\vspace{0.5mm}
\textbf{Task 1: Identify Functions, Parameters \& Return Types} (6 points)
Identify three independent helper functions suitable for this modular decomposition:
\begin{itemize}[leftmargin=1.8em,itemsep=2.5pt]
  \item \textbf{Function 1:} \texttt{validate\_trip(d, w, t)} $\rightarrow$ Return Type: \rule{2.2cm}{0.4pt} | Purpose: \dotfill
  \item \textbf{Function 2:} \texttt{compute\_fuel\_cost(d, w, p)} $\rightarrow$ Return Type: \rule{2.2cm}{0.4pt} | Purpose: \dotfill
  \item \textbf{Function 3:} \texttt{compute\_driver\_pay(t, d, is\_night)} $\rightarrow$ Return Type: \rule{2.2cm}{0.4pt} | Purpose: \dotfill
\end{itemize}

\vspace{0.5mm}
\textbf{Task 2: Write the Structured Modular Call Hierarchy in Ordered Steps} (7 points)\\
Formulate sequential pseudo-logic showing how \texttt{compute\_total\_reimbursement} coordinates validation, invokes helper functions, applies conditional bonuses, and returns the net total.

\begin{answerbox}[height=4.5cm]{Program-Logic Work Area}
\end{answerbox}

\vspace{0.5mm}
\textbf{Task 3: Hand Trace / Step-by-Step Test Calculation} (6 points)\\
For $d = 300.0\text{ km}$, $w = 2500.0\text{ kg}$, $t = 5.0\text{ hours}$, $\text{is\_night} = \text{True}$, and fuel price $p = \$1.50/\text{L}$, calculate the fuel cost, driver pay, and final total reimbursement.

\begin{answerbox}[height=3.0cm]{Calculation Box for Task 3}
\end{answerbox}

\vspace{0.5mm}
\textbf{Task 4: Explain Two Software Engineering Design Decisions} (6 points)
\begin{enumerate}[label=\alph*),leftmargin=1.8em,itemsep=2pt]
  \item Why must functions return calculated values rather than directly calling \texttt{print()} inside helper functions?\\
  \rule{\linewidth}{0.4pt}
  \item Why is passing data via explicit parameters and return values superior to relying on shared global variables?\\
  \rule{\linewidth}{0.4pt}
\end{enumerate}

\newpage

% ==============================================================================
% PAGE 5: SECTION D, PROBLEM 1
% ==============================================================================

\section*{Section D: Applied Programming in Python \hfill \normalsize\normalfont(30 points)}
\textit{Write complete, syntactically correct Python code. Design well-structured modular functions with explicit parameters and return values. Do \textbf{not} use global variables or third-party libraries. Organize code with clean helper functions and a \texttt{main()} entry point. Suggested time: 18 minutes.}

\subsection*{Problem 1: Modular Retail Inventory \& Tiered Volume Billing System (15 points)}
Write a modular Python program that computes customer order charges with volume discounting, weight-based shipping, and sales tax.

\textbf{Requirements:}
\begin{enumerate}[leftmargin=2em,label=\arabic*.,itemsep=1pt]
  \item Implement \texttt{get\_discount\_rate(quantity)}: Returns a float discount rate based on units purchased:
    0\% if $\text{quantity} < 10$; 5\% if $10 \le \text{quantity} < 50$; 10\% if $50 \le \text{quantity} < 100$; 15\% if $\text{quantity} \ge 100$.
  \item Implement \texttt{compute\_shipping(weight, is\_express)}: Base fee is \$5.00 for $\text{weight} \le 5.0$~kg. For weight exceeding 5.0~kg, add \$1.50 per excess kg. If \texttt{is\_express} is \texttt{True}, add a flat \$12.00 express fee.
  \item Implement \texttt{compute\_order\_total(unit\_price, quantity, weight, is\_express, tax\_rate)}:
    \begin{itemize}[itemsep=0.5pt]
      \item Gross subtotal = $\text{unit\_price} \times \text{quantity}$. Discount amount = $\text{gross} \times \text{get\_discount\_rate(quantity)}$.
      \item Net merchandise = $\text{gross} - \text{discount}$.
      \item Shipping fee: calculated via \texttt{compute\_shipping(weight, is\_express)}. \textbf{Free shipping rule:} if net merchandise $\ge \$200.00$ and \texttt{is\_express} is \texttt{False}, shipping fee is waived (\$0.00).
      \item Tax = $\text{net merchandise} \times \text{tax\_rate}$. Total payable = $\text{net merchandise} + \text{shipping} + \text{tax}$.
      \item Return a tuple: \texttt{(gross, discount, net\_merchandise, shipping, tax, total\_payable)}.
    \end{itemize}
  \item Implement \texttt{main()}: Prompt user for inputs, validate positive price/quantity/weight, call \texttt{compute\_order\_total} with 8\% tax rate (\texttt{0.08}), and display a formatted invoice.
\end{enumerate}

Sample formatted output for \$25.00 unit price, 12 units, 8.0~kg, standard delivery (not express):
\begin{lstlisting}
Gross Subtotal: $300.00
Volume Discount (5%): $15.00
Net Merchandise: $285.00
Shipping Fee (Waived): $0.00
Sales Tax (8%): $22.80
Total Payable: $307.80
\end{lstlisting}

\begin{answerbox}[height=7.8cm]{Python Solution for Problem 1}
\end{answerbox}

\newpage

% ==============================================================================
% PAGE 6: SECTION D, PROBLEM 2
% ==============================================================================

\subsection*{Problem 2: Student Academic Standing \& Merit Scholarship Evaluator (15 points)}
Write a modular Python program to evaluate student academic performance and financial scholarship grants.

\textbf{Requirements:}
\begin{enumerate}[leftmargin=1.8em,label=\arabic*.,itemsep=1.5pt]
  \item Implement \texttt{validate\_student\_record(student\_id, gpa, credits)}: Returns \texttt{True} if \texttt{student\_id} has length 8, starts with \texttt{"NEU"}, ends with 5 digits (\texttt{student\_id[3:].isdigit()}), $0.0 \le \text{gpa} \le 4.0$, and $\text{credits} \ge 0$; otherwise \texttt{False}.
  \item Implement \texttt{determine\_standing(gpa, credits)}: Returns a standing string:
    \texttt{"Dean's List"} if $\text{gpa} \ge 3.60$ and $\text{credits} \ge 15$;
    \texttt{"Good Standing"} if $\text{gpa} \ge 2.50$;
    \texttt{"Academic Warning"} if $2.00 \le \text{gpa} < 2.50$;
    or \texttt{"Academic Probation"} if $\text{gpa} < 2.00$.
  \item Implement \texttt{calculate\_scholarship(gpa, credits, is\_financial\_aid)}:
    \begin{itemize}[itemsep=0.5pt]
      \item Merit tier: \$1200.00 if $\text{gpa} \ge 3.80$; \$800.00 if $3.50 \le \text{gpa} < 3.80$; \$400.00 if $3.20 \le \text{gpa} < 3.50$; else \$0.00.
      \item Need supplement: If \texttt{is\_financial\_aid} is \texttt{True} and $\text{gpa} \ge 2.50$, add \$350.00.
      \item Full-time course grant: If $\text{credits} \ge 15$ and initial scholarship $> \$0.00$, add \$150.00.
      \item Return the total calculated scholarship amount as a \texttt{float}.
    \end{itemize}
  \item Implement \texttt{main()}: Prompt user for student details, validate via \texttt{validate\_student\_record}, and display a clean academic report card enclosed with border lines (\texttt{"=" * 45}).
\end{enumerate}

Sample output for \texttt{NEU54321}, GPA \texttt{3.85}, \texttt{16} credits, financial aid \texttt{"yes"}:
\begin{lstlisting}
=============================================
ACADEMIC STANDING & SCHOLARSHIP REPORT
---------------------------------------------
Student ID: NEU54321 | GPA: 3.85 | Credits: 16
Academic Standing: Dean's List
Merit & Need Scholarship: $1700.00
=============================================
\end{lstlisting}

\begin{answerbox}[height=8.5cm]{Python Solution for Problem 2}
\end{answerbox}

\vfill
\begin{center}
  \rule{0.5\textwidth}{0.4pt}\\[1mm]
  \textbf{\color{primaryblue}--- END OF EXAM PAPER ---}
\end{center}

\end{document}
